% Vietnamese translation for OI Public Library
% Source-File: IOI中国国家候选队论文/2005/朱泽园/朱泽园.ppt
\documentclass[11pt]{article}
\usepackage{oipl-vn}

\title{Bản trình chiếu: Trở về điểm xuất phát}
\author{Zhu Zeyuan}
\date{2005}

\begin{document}
\maketitle

\origtitle{Slide companion for back-to-the-starting-point thinking}
\sourcepath{IOI China candidate team papers / 2005 / Zhu Zeyuan / Zhu Zeyuan.ppt}

\section{Phạm vi bản dịch}

Bản trình chiếu đi cùng bài viết về tư duy ``trở về điểm xuất phát''. Phần
chữ đọc được có các ví dụ \textit{USACO Shaping Regions} và
\textit{BalticOI 2004 1-3 Sequence}, cùng các độ phức tạp từ \(O(n^3)\) xuống
\(O(n\log n)\).

\section{Tư tưởng chính}

Khi một hướng giải ngày càng rối, ta quay lại mô hình ban đầu và hỏi lại:
đáp án thật sự phụ thuộc vào đại lượng nào? Việc bỏ một mô hình trung gian
không phù hợp có thể mở ra cách nhìn mới đơn giản hơn.

\section{Ví dụ dãy đơn điệu}

Trong \textit{1-3 Sequence}, cần chọn dãy không giảm
\[
  z_1\le z_2\le \cdots \le z_N
\]
để tối thiểu hóa
\[
  |t_1-z_1|+|t_2-z_2|+\cdots+|t_N-z_N|.
\]
Các slide cho thấy từ cách làm chậm \(O(n^2\log n)\) có thể quay lại cấu trúc
đoạn và giá trị trung vị để đạt \(O(n\log n)\).

\section{Kết luận}

Điểm mạnh của tư duy này là dám phá bỏ giả định cài đặt ban đầu. Bản trình
chiếu dùng các bài hình học và dãy số để minh họa quá trình tìm lại mô hình
cốt lõi.

\end{document}
